\documentclass[10pt,a4paper]{amsart}
\usepackage[margin=19mm]{geometry}
\usepackage{amsmath,amssymb,mathtools}
\usepackage[scr=rsfs,cal=euler]{mathalfa}
\usepackage{xfrac}
\usepackage[unicode,hidelinks,bookmarks=false]{hyperref}
\newcommand{\ie}{i.e.\ }
\newcommand{\cf}{cf.\ }
\newcommand{\resp}{respectively\ }
\newcommand{\Subsection}{\subsection}
\providecommand{\nobreakdash}{\nobreak\hbox{-}}
\setlength{\emergencystretch}{2em}
\multlinegap0pt
\usepackage{bm}
\usepackage{bbm}
\usepackage{dsfont}
\usepackage{xspace}
\usepackage{cancel}
\usepackage{mathtools}
\usepackage{amsthm}
\makeatletter
\@ifundefined{theorem}{\newtheorem{theorem}{Theorem}}{}
\@ifundefined{lemma}{\newtheorem{lemma}[theorem]{Lemma}}{}
\makeatother
\makeatletter
\expandafter\ifx\csname proof\endcsname\relax
\newenvironment{proof}[1][Proof]{\noindent\textbf{#1.} }{\ \rule{0.5em}{0.5em}}
\fi
\makeatother
\title[The major index (maj) and its Sch\"utzenberger dual]{The major index (maj) and its Sch\"utzenberger dual}
\author{Oleg Ogievetsky, Senya Shlosman}
\date{Source: arXiv:2502.02262v2 (CC BY 4.0)}
\begin{document}
\maketitle

\noindent\emph{Source and licence.} The major index (maj) and its Sch\"utzenberger dual. Version arXiv:2502.02262v2 (CC BY 4.0), distributed under the Creative Commons Attribution 4.0 International licence, as recorded in the arXiv licence metadata for this exact version and shown on the abstract page https://arxiv.org/abs/2502.02262v2. Full rights notice: licenses/arxiv-2502.02262-CC-BY-4.0.txt.\par\smallskip

\noindent\textit{Editorial scope.} Counted unit: the Lemma whose source label is \texttt{None} in the article (source0.tex lines 1153--1157, source label \texttt{None}), delivered together with the article's own complete original proof of it (source0.tex lines 1159--1325, one \texttt{proof} environment of 167 lines). No mathematics is written, repaired or re-derived: statement and proof are the article's own text line for line. Required context retained uncounted: none. Every definition, display and label of those lines is retained unchanged. Omitted from this package and not redistributed: 1--1152, 1158--1158, 1326--1386. Nothing inside a retained excerpt is omitted or altered: every retained body is delivered exactly as the article's lines are written, so no omission receipt of any kind accompanies this package. Citations: the retained excerpts carry no citation command at all, so no bibliography entry of the article is transcribed and no reference is redistributed. Reference adaptation: none was needed and none was made; every reference inside the delivered unit resolves inside it, so no unresolved reference or citation is produced. Printed slips kept literal: no printed word, symbol, hypothesis or number of the counted statement or of its proof was altered. Layout and class: the article's own class is \texttt{article} with packages including graphicx, amsmath, amsfonts, amssymb; the wrapper is the lane's pinned \texttt{amsart} editorial wrapper at 10pt on A4, which restates the theorem-like environments the retained text needs and adds the article's own macro definitions that the retained text needs (none), with extra wrapper packages (bm, bbm, dsfont, xspace, cancel, mathtools). The delivered numbering is the unit's own. Assets: no retained span contains a figure environment, a graphics-inclusion command, a PDF-page inclusion, a table or a TikZ picture; no image file is copied into this package, the package declares no asset dependency and no image file is redistributed with it. Licence: version arXiv:2502.02262v2 of this article is distributed under the Creative Commons Attribution 4.0 International licence, as recorded in the arXiv licence metadata for this exact version and shown on the abstract page https://arxiv.org/abs/2502.02262v2. A full rights notice is delivered at licenses/arxiv-2502.02262-CC-BY-4.0.txt. Characters: every retained excerpt is pure ASCII, so the delivered engine, XeLaTeX with its default Latin Modern text font, reports no missing character.
\begin{lemma}
Let $Q^{\prime}$ be a SYT (may be with a hole), and $Q^{\prime\prime}$ be the
SYT obtained from $Q^{\prime}$ by a single jdt move. Then $\mathsf{Des}\left(
Q^{\prime}\right)  =\mathsf{Des}\left(  Q^{\prime\prime}\right)  .$
\end{lemma}

\begin{proof}
Depending on the shape of the SYT $Q^{\prime}$ and the location of the cell
$\ast$ in it, there are four different cases of the initial configuration of
the jdt step:

$%
\begin{tabular}
[c]{|l|l|l|}\hline
$\cdot$ & $\cdot$ & $\cdot$\\\hline
$\cdot$ & $\ast$ & $s$\\\hline
$\cdot$ &  & \\\hline
\end{tabular}
,\
\begin{tabular}
[c]{|l|l|l|}\hline
$\cdot$ & $\cdot$ & $\cdot$\\\hline
$\cdot$ & $\ast$ & $s$\\\hline
$\cdot$ & $t$ & $\cdot$\\\hline
\end{tabular}
,\
\begin{tabular}
[c]{|l|l|l|}\hline
$\cdot$ & $\cdot$ & $\cdot$\\\hline
$\cdot$ & $\ast$ & \\\hline
$\cdot$ & $s$ & \\\hline
\end{tabular}
,\ $and $%
\begin{tabular}
[c]{|l|l|l|}\hline
$\cdot$ & $\cdot$ & $\cdot$\\\hline
$\cdot$ & $\ast$ & $t$\\\hline
$\cdot$ & $s$ & $\cdot$\\\hline
\end{tabular}
.$ Here $s<t$ are integers, and in the first case the cell $\left(
2,3\right)  $ is not part of $Q^{\prime},$ while in the third case the cell
$\left(  3,2\right)  $ is not part of it.

\textbf{1. }In the first case the jdt step is

$%
\begin{tabular}
[c]{|l|l|l|}\hline
$\cdot$ & $\cdot$ & $\cdot$\\\hline
$\cdot$ & $\ast$ & $s$\\\hline
$\cdot$ &  & \\\hline
\end{tabular}
\rightarrow%
\begin{tabular}
[c]{|l|l|l|}\hline
$\cdot$ & $\cdot$ & $\cdot$\\\hline
$\cdot$ & $s$ & $\ast$\\\hline
$\cdot$ &  & \\\hline
\end{tabular}
,$ and since no cell moves vertically, the set $\mathsf{Des}$ does not change.

\textbf{2.} The same applies to the second case, where we have

$%
\begin{tabular}
[c]{|l|l|l|}\hline
$\cdot$ & $\cdot$ & $\cdot$\\\hline
$\cdot$ & $\ast$ & $s$\\\hline
$\cdot$ & $t$ & $\cdot$\\\hline
\end{tabular}
$ $\rightarrow%
\begin{tabular}
[c]{|l|l|l|}\hline
$\cdot$ & $\cdot$ & $\cdot$\\\hline
$\cdot$ & $s$ & $\ast$\\\hline
$\cdot$ & $t$ & $\cdot$\\\hline
\end{tabular}
.$

\textbf{3. }Consider the jdt step

$%
\begin{tabular}
[c]{|l|l|l|}\hline
$\cdot$ & $\cdot$ & $\cdot$\\\hline
$\cdot$ & $\ast$ & \\\hline
$\cdot$ & $s$ & \\\hline
\end{tabular}
\rightarrow%
\begin{tabular}
[c]{|l|l|l|}\hline
$\cdot$ & $\cdot$ & $\cdot$\\\hline
$\cdot$ & $s$ & \\\hline
$\cdot$ & $\ast$ & \\\hline
\end{tabular}
.$ Note that the only two values which we have to worry about are $s$ and
$s-1.$

If $s$ is in $\mathsf{Des}\left(  Q^{\prime}\right)  ,$ then $s$ is in
$\mathsf{Des}\left(  Q^{\prime\prime}\right)  ,$ since $s$ is ascending, and
therefore remains a descent.

Note that the cell $Q^{\prime}\left(  s+1\right)  $ does not belong to the row
of the cell $Q^{\prime}\left(  s\right)  $ -- indeed, it cannot be in this row
and to the left of $Q^{\prime}\left(  s\right)  $ since $s+1>s,$ while
$Q^{\prime}\left(  s\right)  $ is the rightmost cell of that row. Therefore if
$s\notin\mathsf{Des}\left(  Q^{\prime}\right)  ,$ then $s\notin\mathsf{Des}%
\left(  Q^{\prime\prime}\right)  .$

If the cell $Q^{\prime}\left(  s-1\right)  $ is not a descent, then the upward
move of the cell $Q^{\prime}\left(  s\right)  $ cannot make it to be one.

If the cell $Q^{\prime}\left(  s-1\right)  $ is two or more rows higher than
the cell $Q^{\prime}\left(  s\right)  ,$ then the value $s-1$ is a descent for
$Q^{\prime}$ as well as for $Q^{\prime\prime}.$

Finally, the cell $Q^{\prime}\left(  s-1\right)  $ cannot be one row higher
than $Q^{\prime}\left(  s\right)  $ -- indeed, all the values to the left of
$\ast$ in $Q^{\prime}$ are $\leq s-2,$ since for the fragment $%
\begin{tabular}
[c]{|l|l|l|}\hline
$\cdot$ & $\cdot$ & $\cdot$\\\hline
$x$ & $\ast$ & \\\hline
$y$ & $s$ & \\\hline
\end{tabular}
$ of $Q^{\prime}$ we have $x<y<s.$ So $s-1\in\mathsf{Des}\left(  Q^{\prime
}\right)  \ $iff $s-1\in\mathsf{Des}\left(  Q^{\prime\prime}\right)  .$

\textbf{4. }The remaining jdt step to consider is

$%
\begin{tabular}
[c]{|l|l|l|}\hline
$\cdot$ & $\cdot$ & $\cdot$\\\hline
$\cdot$ & $\ast$ & $t$\\\hline
$\cdot$ & $s$ & $\cdot$\\\hline
\end{tabular}
\rightarrow%
\begin{tabular}
[c]{|l|l|l|}\hline
$\cdot$ & $\cdot$ & $\cdot$\\\hline
$\cdot$ & $s$ & $t$\\\hline
$\cdot$ & $\ast$ & $\cdot$\\\hline
\end{tabular}
.$

As in \textbf{3}, if $s$ is in $\mathsf{Des}\left(  Q^{\prime}\right)  ,$ then
$s$ is in $\mathsf{Des}\left(  Q^{\prime\prime}\right)  ,$ since $s$ has ascended.

Again, as in \textbf{3}, the cell $Q^{\prime}\left(  s+1\right)  $ does not
belong to the row of the cell $Q^{\prime}\left(  s\right)  $ -- it cannot be
in this row and to the left of $Q^{\prime}\left(  s\right)  $ since $s+1>s,$
while every entry $x$ in that row to the right of $Q^{\prime}\left(  s\right)
$ satisfy $t<x,$ so $x\geq s+2,$ thus if $s\notin\mathsf{Des}\left(
Q^{\prime}\right)  $ then $s\notin\mathsf{Des}\left(  Q^{\prime\prime}\right)
.$

If the cell $Q^{\prime}\left(  s-1\right)  $ is not a descent, then the upward
move of the cell $Q^{\prime}\left(  s\right)  $ cannot make it to be one.

If the cell $Q^{\prime}\left(  s-1\right)  $ is two or more rows higher than
the cell $Q^{\prime}\left(  s\right)  ,$ then $s-1\in\mathsf{Des}\left(
Q^{\prime}\right)  $ and $s-1\in\mathsf{Des}\left(  Q^{\prime\prime}\right)
.$ Finally, the cell $Q^{\prime}\left(  s-1\right)  $ cannot be one row higher
than $Q^{\prime}\left(  s\right)  $ -- indeed, $t>s,$ while for the fragment $%
\begin{tabular}
[c]{|l|l|l|}\hline
$\cdot$ & $\cdot$ & $\cdot$\\\hline
$x$ & $\ast$ & $t$\\\hline
$y$ & $s$ & $\cdot$\\\hline
\end{tabular}
$ of $Q^{\prime}$ we have $x<y<s.$
\end{proof}

\end{document}
